\chapter{Method, sources and limits}\label{ch:methods}
\section{Choose a method for the question}
Explain what your method allows you to establish. A doctrinal inquiry may
reconstruct the legal conditions of a power or duty. A comparison may test
whether different institutions perform a sufficiently similar function to
justify analysis together. An empirical inquiry needs a defensible design for
observing practices and interpreting the observations. A thesis combining these
approaches should explain how their findings connect.

Comparative method should be selected with the research question and purpose
in view. The example reference below points to scholarship addressing that
relationship.\lawcite[1]{vanhoecke2015}

\section{Define the source corpus}
Identify the types of primary material included, the databases and languages
searched, relevant date limits and reasons for exclusion. Explain how amended
legislation, subsequent decisions, official guidance and persuasive sources
are distinguished. Record an instrument's version, commencement position and
territorial scope where these affect the argument.

\begin{table}[htbp]
\centering\small
\caption{Questions to resolve when selecting legal materials}\label{tab:corpus}
\begin{tabularx}{\textwidth}{@{}p{0.25\textwidth}X@{}}
\toprule
Decision & Record needed \\
\midrule
Authority & Issuing body, jurisdiction, legal status and relevant hierarchy \\
Time & Decision date or legislative version, cut-off and later treatment \\
Language & Original text, translation source and translator's role \\
Selection & Search route, inclusion criterion and reason for exclusion \\
Verification & Exact proposition, pinpoint, reviewer and verification date \\
\bottomrule
\end{tabularx}
\end{table}

\section{Interpretation and disagreement}
Separate what an authority says from your inference about its effect. Identify
the interpretive method, institutional context and competing reading. A treaty
provision can be cited with an article pinpoint.\lawcite[art~31(1)]{vclt1969}
This demonstration is a reference to an identified provision; applying that
provision to a particular controversy requires substantive analysis.

\section{Research integrity and limitations}
Keep records sufficient to reconstruct important research decisions. Explain
how sources, assistance and limitations are reported. The European Code of
Conduct provides a general integrity framework to consult alongside applicable
institutional requirements.\lawcite[ss~2.3 and 2.5]{allea2023}

\LawPrompt{State material limits concerning language, access, selection,
generalisability and interpretation. Explain how each affects the conclusions.}
